\subsection{线性映射的紧群}

设 E 是 R、C 或 H 上的有限维向量空间。则 End(E) 是 R 上的有限维代数，而 End(E) 上的典范拓扑（§1, No.~10）就是紧收敛拓扑。群 \(\operatorname{Aut}(E) = \mathbf{GL}(E)\) 是 End(E) 的开子集，因而是局部紧群。设 \((e_{1}, e_{2}, \ldots, e_{n})\) 是 E 的一组基，对于 E 的每个自同态 u，令 \(M(u) = (\alpha_{ij}(u))\) 为 u 关于此基的矩阵；End(E) 的子集 S 在 End(E) 中相对紧，等价于函数 \(\alpha_{ij}(u)\) 在 S 上有界。

命题 1. --- 设 G 是 Aut(E) 的子群。以下三个性质等价：

\begin{enumerate}
\def\labelenumi{(\roman{enumi})}
\item
  G 在 \(\operatorname{End}(\mathbf{E})\) 中相对紧；
\item
  G 在 \(\operatorname{Aut}(\mathbf{E})\) 中相对紧；
\item
  G 保持 E 上的一个非退化 \(^{1}\) 正厄米形式不变。(iii) \(\Rightarrow\) (i)：假设 G 保持非退化正厄米形式 \(\Psi\) 不变。取 \((e_{1},\ldots,e_{n})\) 为 \(\Psi\) 的一组标准正交基（Alg., Ch. IX, §6, No.~1, Cor. 1 of Th. 1）。对于每个 \(u\in G\)，令 \((u_{ij})\) 为 u 关于 \((e_{i})\) 的矩阵。对于任意 j，有 \(\sum_{i=1}^{n}|u_{ij}|^{2}=1\)，故对于所有 i 和 j 都有 \(|u_{ij}|\leqslant1\)，这就证明了 (i)。
\item
  \(\Rightarrow\) (ii)：这由 GT, X, §3, No.~5, Th. 4 的 Cor. 得出，同时注意到 \(\operatorname{End}(\mathbf{E})\) 的拓扑是紧收敛拓扑。
\item
  \(\Rightarrow\) (iii)：假设 \(\overline{\mathbf{G}}\) 是 \(\mathbf{G}\) 在 \(\operatorname{Aut}(\mathbf{E})\) 中的闭包且是紧的。设 \(\Phi\) 是 \(\mathbf{E}\) 上的非退化正厄米形式。若标量域是 \(\mathbf{R}\) 或 \(\mathbf{C}\)，则由 \(\Phi\) 使 \(\mathbf{E}\) 成为有限维 Hilbert 空间，性质 (iii) 将由下述引理得到：
\end{enumerate}

引理 1. --- 设 F 是 Hilbert 空间，K 是紧群，且 \(s \mapsto U(s)\) 是 K 到 \(\mathcal{L}(\mathrm{F};\mathrm{F})\) 的可逆元群中的表示，并且关于逐点收敛拓扑连续。则 F 上存在非退化正厄米形式 \(\varphi\)，使得

\[
\varphi (U (s) x, U (s) y) = \varphi (x, y)
\]

对于所有 \(s \in \mathbf{K}\)、\(x \in \mathbf{F}\) 和 \(y \in \mathbf{F}\) 成立，并且 \(\mathbf{F}\) 上由 \(\varphi\) 定义的拓扑向量空间结构（TVS, V, §1, No.~3）与 \(\mathbf{F}\) 的原有结构相同。

设 \(\alpha\) 是 \(\mathbf{K}\) 上的 Haar 测度。对于任意 \(x, y\)，在 \(\mathbf{F}\) 中，映射 \(s \mapsto (U(s)x|U(s)y)\) 连续。置

\[
\varphi (x, y) = \int (U (s) x | U (s) y) d \alpha (s).
\]

显然，\(\varphi(x,y)\) 是 \(\mathbf{F}\) 上的半双线性形式。由于自同态集 \(U(s)\) 在 \(\mathcal{L}_s(\mathrm{F};\mathrm{F})\) 中紧，存在常数 \(\mathbf{M}\)，使得 \(\| U(s)\| \leqslant \mathbf{M}\) 对所有 \(s\in \mathbf{K}\) 成立。因此对于每个 \(x\in \mathbf{F}\)，有

\[
\mathrm{M} ^ {- 1} \| x \| \leqslant \| U (s) x \| \leqslant \mathrm{M} \| x \|,
\]

从而有不等式

\[
\mathbf {M} ^ {- 2} \alpha (\mathbf {K}) \| x \| ^ {2} \leqslant \varphi (x, x) \leqslant \mathbf {M} ^ {2} \alpha (\mathbf {K}) \| x \| ^ {2},
\]

这表明 \(\varphi\) 是正且非退化的，并且范数 \(\varphi(x,x)^{1/2}\) 与范数 \(\|x\|\) 等价。最后，对于所有 \(t \in K\)，

\[
\begin{array}{r l} & {\varphi (U (t) x, U (t) y) = \int (U (s t) x | U (s t) y) d \alpha (s)} \\ & {\qquad = \int (U (s) x | U (s) y) d \alpha (s) = \varphi (x, y).} \end{array}
\]

当标量域为 H 时，只需完全类似地处处将函数 \(s \mapsto (U(s)x|U(s)y)\) 换成定义在 G 上、取值于 H 的函数 \(s \mapsto \Phi(sx, sy)\)。命题得证。

注记。--- 设 \(\Phi\) 是 E 上的非退化正厄米形式。酉群 \(\mathbf{U}(\Phi)\) 在 \(\operatorname{Aut}(\mathbf{E})\) 中是闭的，因而是紧的（Prop. 1）。Prop. 1 还表明，\(\operatorname{Aut}(\mathbf{E})\) 的每个紧子群都包含于某个形如 \(\mathbf{U}(\Phi)\) 的子群中。现在若 \(\mathbf{U}(\Phi)\) 包含于 \(\operatorname{Aut}(\mathbf{E})\) 的紧子群 K 中，则可见存在 E 上的非退化正厄米形式 \(\Phi'\)，使得 \(\mathbf{U}(\Phi) \subset \mathbf{K} \subset \mathbf{U}(\Phi')\)；并且容易推出（Exer. 1），\(\Phi\) 与 \(\Phi'\) 成比例，从而 \(\mathbf{U}(\Phi) = K\)。因此，\(\operatorname{Aut}(E)\) 的极大紧子群正是形如 \(\mathbf{U}(\Phi)\) 的子群。

